\documentclass[11pt]{article}
\usepackage[T1]{fontenc}
\usepackage[utf8]{inputenc}
\usepackage{lmodern}
\usepackage{amsmath,amssymb,amsthm}
\usepackage[margin=1in]{geometry}
\usepackage{booktabs}
\usepackage{xcolor}
\usepackage[colorlinks=true,linkcolor=blue,citecolor=blue,urlcolor=blue]{hyperref}
\newtheorem{theorem}{Theorem}
\newtheorem{lemma}{Lemma}
\newtheorem{proposition}{Proposition}
\newcommand{\one}{\mathbf 1}
\newcommand{\EE}{\mathbb E}
\newcommand{\PP}{\mathbb P}
\newcommand{\diag}{\operatorname{diag}}
\newcommand{\CV}{\operatorname{CV}}
\title{Two counterexamples for\\stationary MMPP arrival times}
\author{}
\date{Research note for review\quad---\quad 2 October 2026}
\begin{document}
\maketitle
\begin{abstract}
Conjecture 3.9 of Asanjarani and Nazarathy, published in the ANZIAM
Journal in 2022, asserts that the time-stationary first-arrival time of
an irreducible Markov-modulated Poisson process stochastically dominates
its event-stationary first-arrival time. We exhibit an irreducible
16-state directed-cycle model with rational transition and arrival
rates for which the asserted survival inequality is reversed at time
$12$. A self-contained analytic argument proves a violation exceeding
$5.5\times10^{-6}$, and a separate exact rational matrix-series
certificate verifies the sign. The same rational model also answers negatively He's separate question
whether the time-stationary first-arrival time must have squared
coefficient of variation at least one. This leaves intact the
event-stationary squared-coefficient-of-variation result established
by He. The mathematical verification below does not assert priority;
independent literature and human review remain necessary before any
publication claim.
\end{abstract}

\section{The conjecture and the model}
We use row-vector conventions for probability distributions. An
irreducible finite-state Markov-modulated Poisson process (MMPP) consists
of a continuous-time Markov chain $J$ with generator $Q$, and a Poisson
arrival process whose conditional instantaneous rate is $d_{J_t}$.
Write
\[
 D=\diag(d_0,\ldots,d_{15}),\qquad C=Q-D,
 \qquad \pi Q=0,\quad \pi\one=1.
\]
The mean arrival rate and the event-stationary state distribution are
\begin{equation}\label{eq:palm}
 \overline\lambda=\pi D\one,\qquad
 \alpha=\frac{\pi D}{\overline\lambda}.
\end{equation}
For an MMPP, arrivals do not change the state of $J$, so $\alpha$ is
also the state distribution immediately after an arrival in the
stationary arrival process. If $T^\beta$ is the time to the next arrival
when $J_0$ has distribution $\beta$, its survival function is
\begin{equation}\label{eq:survival}
 F_\beta(t)=\PP(T^\beta>t)=\beta e^{tC}\one
 =\EE_\beta\exp\left(-\int_0^t d_{J_s}\,ds\right).
\end{equation}
Conjecture 3.9 of \cite{AN} states $T^\pi\geq_{\rm st}T^\alpha$ for every
irreducible MMPP, equivalently
\begin{equation}\label{eq:conjecture}
 F_\pi(t)\geq F_\alpha(t)\qquad(t\geq0).
\end{equation}
Equation \eqref{eq:palm} and the survival inequality are equations
(2.1) and (2.9), respectively, in that paper. No reversibility of $Q$
is assumed in its conjecture.

Index the states modulo $16$. Set
\begin{equation}\label{eq:Q}
 Q_{i,i}=-1,\qquad Q_{i,i+1}=1,
 \qquad Q_{i,j}=0\quad(j\notin\{i,i+1\}).
\end{equation}
This generator is irreducible, and its stationary distribution is
$\pi_i=1/16$. Define the integer vector
\begin{equation}\label{eq:s}
\begin{split}
 (s_0,\ldots,s_{15})={}&(1000,924,707,383,0,-383,-707,-924,\\
 &-1000,-924,-707,-383,0,383,707,924)
\end{split}
\end{equation}
and choose the arrival rates
\begin{equation}\label{eq:rates}
 d_i=\frac{10000+s_i}{10^6}
 =\lambda+\varepsilon g_i,
 \qquad \lambda=\frac1{100},\quad
 \varepsilon=\frac1{1000},\quad g_i=\frac{s_i}{1000}.
\end{equation}
Since $\sum_i s_i=0$, we have $\overline\lambda=\lambda$ and
\[
 \alpha_i=\frac{10000+s_i}{160000}.
\]
Every arrival rate lies in $[9/1000,11/1000]$, so the model has strictly
positive rates. In particular, both $\pi$ and $\alpha$ are probability
distributions with strictly positive entries.

\begin{theorem}\label{thm:main}
For the irreducible MMPP \eqref{eq:Q}--\eqref{eq:rates},
\begin{equation}\label{eq:violation}
 F_\pi(12)-F_\alpha(12)<-5.544\times10^{-6}.
\end{equation}
Thus this model is a counterexample to \eqref{eq:conjecture}.
\end{theorem}

\section{Survival, hazard, and perturbation identities}
We first record identities valid for any finite-state stationary MMPP
with positive arrival rates. Since $\pi C=-\pi D$, differentiating
\eqref{eq:survival} gives
\begin{equation}\label{eq:derivative}
 F_\pi'(t)=\pi C e^{tC}\one
 =-\pi D e^{tC}\one
 =-\overline\lambda F_\alpha(t).
\end{equation}
This uses the stationary identity $\pi Q=0$, not reversibility. The
hazard rate from the time-stationary distribution therefore satisfies
\begin{equation}\label{eq:hazard}
 h_\pi(t):=-\frac{F_\pi'(t)}{F_\pi(t)}
 =\overline\lambda\frac{F_\alpha(t)}{F_\pi(t)}.
\end{equation}
Consequently, the asserted order \eqref{eq:conjecture} is equivalent to
$h_\pi(t)\leq\overline\lambda$ for all $t\geq0$.

For the explicit model, take $J_0\sim\pi$ and put
\begin{equation}\label{eq:H}
 A_t=\int_0^t g(J_s)\,ds,\qquad
 H(t)=\EE_\pi e^{-\varepsilon A_t}.
\end{equation}
Then $F_\pi(t)=e^{-\lambda t}H(t)$, and
\eqref{eq:derivative} yields the exact identity
\begin{equation}\label{eq:difference}
 F_\pi(t)-F_\alpha(t)=\frac{e^{-\lambda t}}{\lambda}H'(t).
\end{equation}
The mean of $g(J_t)$ is zero and $|g(J_t)|\leq1$.

\begin{lemma}\label{lem:expansion}
For every $t\geq0$,
\begin{equation}\label{eq:expansion}
 H'(t)=\varepsilon^2 I_g(t)+R(t),\qquad
 I_g(t)=\int_0^t\EE_\pi[g(J_0)g(J_s)]\,ds,
\end{equation}
where
\begin{equation}\label{eq:remainder}
 |R(t)|\leq\frac{\varepsilon^3t^2}{2}e^{\varepsilon t}.
\end{equation}
\end{lemma}
\begin{proof}
For each fixed $t$, the chain has probability zero of jumping at
exactly $t$. Since $|A_{t+h}-A_t|\leq|h|$, dominated differentiation
in \eqref{eq:H} is justified and gives
\[
 H'(t)=-\varepsilon\EE_\pi[g(J_t)e^{-\varepsilon A_t}].
\]
The elementary Taylor formula
\[
 e^{-x}=1-x+x^2\int_0^1(1-u)e^{-ux}\,du
\]
implies $|e^{-x}-1+x|\leq x^2e^{|x|}/2$ for real $x$.
Using $\EE_\pi g(J_t)=0$ and $|A_t|\leq t$ therefore gives
\[
 H'(t)=\varepsilon^2\EE_\pi[g(J_t)A_t]+R(t)
\]
with the bound \eqref{eq:remainder}. By stationarity and Fubini,
\[
 \EE_\pi[g(J_t)A_t]
 =\int_0^t\EE_\pi[g(J_s)g(J_t)]\,ds
 =\int_0^t\EE_\pi[g(J_0)g(J_{t-s})]\,ds=I_g(t).
\]
The equality uses scalar multiplication and a forward time shift;
no time reversal of the chain is needed.
\end{proof}

\section{A self-contained analytic certificate}
Let $\theta=\pi/8$ and $f_i=\cos(i\theta)$. For the directed-cycle
chain \eqref{eq:Q}, one can write
$J_t=J_0+N_t\pmod{16}$, where $N_t$ is a rate-one Poisson process
independent of $J_0$. Averaging over $J_0\sim\pi$ gives
\[
 \frac1{16}\sum_{i=0}^{15}\cos(i\theta)\cos((i+n)\theta)
 =\frac12\cos(n\theta).
\]
The Poisson generating function consequently yields
\begin{equation}\label{eq:covariance}
 \EE_\pi[f(J_0)f(J_t)]
 =\frac12 e^{-at}\cos(bt),\qquad
 a=1-\cos\theta,\quad b=\sin\theta.
\end{equation}
Here $a^2+b^2=2a$. Integrating \eqref{eq:covariance} gives
\begin{equation}\label{eq:integral}
 I_f(t):=\int_0^t\EE_\pi[f(J_0)f(J_s)]\,ds
 =\frac{a+e^{-at}[-a\cos(bt)+b\sin(bt)]}{4a}.
\end{equation}

All terminating decimals below denote exact rational numbers. The
standard bounds $3.141<\pi<3.142$ and the half-angle formulas give
\begin{equation}\label{eq:ab}
 .075<a<.077,\qquad .382<b<.384.
\end{equation}
For clarity, these radical estimates need no numerical approximation:
$1.414^2<2<1.415^2$, and
\[
 \cos^2\theta=\frac{2+\sqrt2}{4},\qquad
 \sin^2\theta=\frac{2-\sqrt2}{4}.
\]
Comparison with $.923^2,.925^2$ proves the bounds on $a$;
comparison with $.382^2,.384^2$ proves the bounds on $b$.

Set $\delta=3\pi/2-12b$. Equation \eqref{eq:ab} gives
\[
 .1035<\delta<.129<.13.
\]
Thus $\cos(12b)=-\sin\delta\geq-.13$ and
\[
 \sin(12b)=-\cos\delta
 \leq-\left(1-\frac{\delta^2}{2}\right)<-.99.
\]
In particular,
\begin{equation}\label{eq:X}
 X:=-a\cos(12b)+b\sin(12b)
 <.077\cdot.13-.382\cdot.99=-.36817.
\end{equation}
Since $12a<.924<1$ and $e<3$, we have $e^{-12a}>1/3$.
The fact that $X$ is negative matters when multiplying inequalities:
\[
 a+e^{-12a}X<.077-\frac{.36817}{3}
 =-\frac{.13717}{3}.
\]
Since $4a<.308$, equation \eqref{eq:integral} now gives
\begin{equation}\label{eq:If}
 I_f(12)<-\frac{.13717}{.924}<-.148.
\end{equation}
The last comparison is exact: $.148\cdot.924=.136752<.13717$.
The estimate $e<3$ follows, for example, from
$\sum_{n\geq2}1/n!<\sum_{n\geq2}2^{-(n-1)}=1$.

We next transfer this sign to the rational vector $g$.
\begin{equation}\label{eq:rounding}
 \|g-f\|_\infty<\frac1{2000}.
\end{equation}
To verify every entry, it suffices by symmetry to consider
$1,\cos\theta,\cos(\pi/4),\sin\theta,0$.
The exact square comparisons $1.4142^2<2<1.4143^2$ give
\[
 .9235<\cos\theta<.9245,\qquad
 .7065<\cos(\pi/4)<.7075,\qquad
 .3825<\sin\theta<.3835.
\]
These prove \eqref{eq:rounding} for the entries $.924,.707,.383$
and hence for all $16$ entries.
For any two times $u,v$, since $|f|,|g|\leq1$,
\[
 |g(J_u)g(J_v)-f(J_u)f(J_v)|
 \leq |g(J_u)-f(J_u)|\,|g(J_v)|
      +|f(J_u)|\,|g(J_v)-f(J_v)|<\frac1{1000}.
\]
It follows that $|I_g(12)-I_f(12)|<.012$. In conjunction with
\eqref{eq:If}, this proves
\begin{equation}\label{eq:Ig}
 I_g(12)<-.136.
\end{equation}

Finally, the remainder in Lemma \ref{lem:expansion} satisfies
\[
 \frac{|R(12)|}{\varepsilon^2}
 \leq .072e^{.012}
 \leq\frac{.072}{1-.012}
 =\frac{18}{247}<.073.
\]
Here $e^x\leq(1-x)^{-1}$ for $0\leq x<1$ follows by comparing their
power series term by term. Therefore
\begin{equation}\label{eq:Hnegative}
 \frac{H'(12)}{\varepsilon^2}<-.136+.073=-.063.
\end{equation}
Using $e^{-.12}>1-.12=.88$ in \eqref{eq:difference} gives
\[
 F_\pi(12)-F_\alpha(12)
 <-.063\cdot10^{-6}\cdot\frac{.88}{.01}
 =-5.544\times10^{-6}.
\]
This completes the proof of Theorem \ref{thm:main} without relying
on floating-point evaluation of a matrix exponential.

\section{An independent exact rational matrix certificate}
The same sign can be verified directly from the rational model.
Set $w=\pi-\alpha$, $r=1011/1000$, and
\[
 B=12(C+rI).
\]
The only nonzero entries of $B$ are
\[
 B_{i,i}=12\left(\frac{11}{1000}-d_i\right)\geq0,
 \qquad B_{i,i+1}=12.
\]
Hence $B$ is entrywise nonnegative and
\[
 L:=\|B\|_\infty=\frac{1503}{125}=12.024,
 \qquad\|w\|_1=\frac{1257}{20000}.
\]
Define the exact rational numbers
\begin{equation}\label{eq:series}
 S_{80}=\sum_{k=0}^{80}\frac{wB^k\one}{k!},\qquad
 E_{80}=\|w\|_1\frac{L^{81}}{81!}\frac1{1-L/82}.
\end{equation}
The matrix norm inequality gives
$|wB^k\one|\leq\|w\|_1L^k$. For $k\geq81$, consecutive scalar
terms $L^k/k!$ have ratio at most $L/82<1$. Summing this geometric
majorant proves
\begin{equation}\label{eq:tail}
 |we^B\one-S_{80}|\leq E_{80}.
\end{equation}
No cancellation assumption is used in the tail bound.

The companion file \texttt{mmpp\_conjecture39\_check.py} evaluates
\eqref{eq:series} using Python's standard-library rational type
\texttt{fractions.Fraction}. It verifies the three exact inequalities
\begin{equation}\label{eq:checks}
 E_{80}<10^{-31},\qquad -4<S_{80}-E_{80},\qquad
 S_{80}+E_{80}<-\frac{39}{10}.
\end{equation}
Its sparse recurrence can also be reproduced independently. Start with
the row vector $u^{(0)}=w$ and, for $k=1,\ldots,80$, set
\[
 u^{(k)}_j=
 \frac{u^{(k-1)}_jB_{j,j}+12u^{(k-1)}_{j-1}}{k}
 \qquad(j\pmod{16}).
\]
Then $u^{(k)}=wB^k/k!$ and
$S_{80}=\sum_{k=0}^{80}\sum_j u^{(k)}_j$. Every operation is rational.
For orientation, the exact rational quantities have decimal displays
\[
 S_{80}\approx-3.9292649153377824801,
 \qquad E_{80}\approx3.8718957312181850\times10^{-35}.
\]
These decimal displays are not used to certify \eqref{eq:checks}.
Combining \eqref{eq:tail} and \eqref{eq:checks} gives
\[
 -4<we^B\one<-\frac{39}{10},\qquad
 F_\pi(12)-F_\alpha(12)=e^{-12r}we^B\one<0.
\]
For a rational negative bound from this certificate alone, $12r<13$
and $e<3$ imply
\[
 F_\pi(12)-F_\alpha(12)
 <-\frac{39}{10\cdot3^{13}}
 =-\frac{13}{5314410}<-2.4\times10^{-6}.
\]
An ordinary numerical evaluation gives
\[
 F_\pi(12)-F_\alpha(12)\approx-2.1156832900682\times10^{-5}.
\]
The analytic and exact rational certificates suffice independently of
this approximation.

\section{Time-stationary squared coefficient of variation}
Section 4.3 of \cite{He} also asks whether the time-stationary
first-arrival time always has squared coefficient of variation at least
one. The same rational model provides a negative answer.

\begin{proposition}\label{prop:scv}
For the MMPP \eqref{eq:Q}--\eqref{eq:rates},
\[
 0<\CV(T^\pi)^2<\frac{199999}{200000}<1.
\]
\end{proposition}
\begin{proof}
Write $R=(D-Q)^{-1}$. Since all arrival rates are at least $9/1000$,
\eqref{eq:survival} implies that $e^{tC}\one\leq e^{-9t/1000}\one$.
Thus the matrix integral $\int_0^\infty e^{tC}\,dt$ converges, equals
$R$, and has nonnegative entries. Integration of the survival function
and the semigroup identity give the usual phase-type moment formulas
\[
 \mu:=\EE_\pi T=\pi R\one,\qquad
 m_2:=\EE_\pi T^2=2\pi R^2\one.
\]
For completeness, the second formula follows from
$R^2=\int_0^\infty\int_0^\infty e^{(s+t)C}\,ds\,dt
=\int_0^\infty t e^{tC}\,dt$.
Let $u=R\one$ and $z=Ru$. Their entries are rational and can be
computed using the following explicit cyclic recurrence. For any
rational right-hand side $b$, the system $(D-Q)x=b$ is
\[
 a_i x_i-x_{i+1}=b_i,\qquad a_i=1+d_i>1,
 \qquad i\pmod{16}.
\]
Put $P=\prod_{i=0}^{15}a_i$. Since $P>1$, the unique solution starts at
\[
 x_0=\frac{\displaystyle\sum_{j=0}^{15}b_j
                    \prod_{k=j+1}^{15}a_k}{P-1},
 \qquad x_{i+1}=a_ix_i-b_i.
\]
Empty products equal one. Iterating the recurrence gives
$x_{16}=Px_0-\sum_j b_j\prod_{k=j+1}^{15}a_k=x_0$, verifying the
cyclic boundary condition. Apply it first to $b=\one$ to obtain $u$,
and then to $b=u$ to obtain $z$.

The companion exact rational verifier evaluates $\mu=\pi u$ and
$m_2=2\pi z$ and checks
\[
 \mu^2<m_2<\frac{399999}{200000}\mu^2.
\]
Dividing by $\mu^2>0$ proves the assertion because
$\CV(T^\pi)^2=m_2/\mu^2-1$. In addition to the cyclic recurrence,
independent exact rational Gaussian elimination gives the same values.
Their approximate decimal displays are
\[
 \mu\approx100.0027991122752187,\qquad
 m_2\approx20001.06092113849143,\qquad
 \CV(T^\pi)^2\approx .9999941263846573.
\]
The strict rational comparisons, rather than these decimal displays,
are the certificate.
\end{proof}

\section{Scope and interpretation}
The sign mechanism is the negative integrated autocovariance of a
nonreversible cycle. For a small modulation of the rate, this quantity
is the leading term in $H'(t)$, and therefore determines which side
of the mean rate the time-stationary hazard occupies. The calculation
above controls all higher-order terms for an explicit rational model.

Theorem \ref{thm:main} disproves the universal survival-order statement
of Conjecture 3.9 in \cite{AN}, and Proposition \ref{prop:scv} answers
negatively the time-stationary variability question in Section 4.3 of
\cite{He}. Neither result refutes the distinct assertion that the
event-stationary interarrival time has squared coefficient of variation
at least one. That assertion is established in \cite{He}; Section 4.3
there treats the stronger stochastic-order question as a further
research question. A survival inequality at every time is stronger
than an inequality between these moments. The directed-cycle example
also makes no claim about reversible modulating chains.

This is a specific published conjecture from 2019/2022, rather than a
famous decades-old open problem. The results and their exact certificates
are presented for review. No claim that this is the first such
counterexample is made; an independent check of subsequent literature
and of the argument is appropriate before submission or announcement.

\begin{thebibliography}{9}
\bibitem{AN}
A.~Asanjarani and Y.~Nazarathy,
\emph{Stationary Markovian Arrival Processes: Results and Open Problems},
The ANZIAM Journal \textbf{64} (2022), 54--68.
\href{https://doi.org/10.1017/S1446181122000013}{doi:10.1017/S1446181122000013}.
Preprint: \href{https://arxiv.org/abs/1905.01736}{arXiv:1905.01736} (2019).

\bibitem{He}
Q.-M.~He,
\emph{Bounds on the mean and squared coefficient of variation of
phase-type distributions},
Journal of Applied Probability \textbf{58} (2021), 880--889.
\href{https://doi.org/10.1017/jpr.2021.18}{doi:10.1017/jpr.2021.18}.
\end{thebibliography}
\end{document}
